\chapter*{머리말}
\addcontentsline{toc}{chapter}{머리말}

현대대수학은 추상적인 정의의 목록이 아니다. 그 출발점에는 방정식, 대칭, 그리고 서로 다른 계산을 하나의 구조로 이해하려는 질문이 있다. 이 책은 그 흐름을 따라 군, 환, 아이디얼과 다항식환을 소개하고, 다음 권에서 체확장과 갈루아 이론으로 나아갈 기반을 마련한다.

\section*{이 책의 독자}

이 책은 다음 독자를 대상으로 한다.

\begin{itemize}
  \item 현대대수학을 처음 배우는 수학과·과학기술계열 학부생
  \item 표준 교재의 압축된 증명 때문에 독학에 어려움을 겪은 학습자
  \item 체론과 갈루아 이론을 공부하기 전에 군론과 환론을 다시 정리하려는 독자
  \item 수학적 글쓰기와 증명 구성 능력을 훈련하려는 대학원 준비생
\end{itemize}

집합과 함수, 정수의 기본 성질, 행렬과 기초 선형대수에 익숙하면 시작할 수 있다. 각 장에서 새로운 선수지식이 필요할 때에는 본문 안에서 가능한 한 다시 설명한다.

\section*{설명의 순서}

많은 교재가 정의를 먼저 제시한 뒤 예제를 붙인다. 이 책은 가능하면 다음 순서를 따른다.

\begin{center}
동기와 질문 $\longrightarrow$ 구체적인 예제 $\longrightarrow$ 정의
$\longrightarrow$ 핵심 정리 $\longrightarrow$ 증명 전략
$\longrightarrow$ 엄밀한 증명 $\longrightarrow$ 연습문제
\end{center}

증명은 정답을 확인하는 장식이 아니라 개념이 작동하는 이유를 보여 주는 본문이다. 따라서 중요한 증명에는 먼저 \textbf{증명 전략}을 제시하고, 그 뒤에 가정을 어디에서 사용하는지 드러나는 완전한 논증을 둔다. 반례가 중요한 개념에는 \textbf{자주 하는 실수} 상자를 두었다.

\section*{권장 학습법}

각 절은 다음과 같이 공부하는 것을 권한다.

\begin{enumerate}
  \item 학습 목표와 핵심 아이디어를 먼저 읽는다.
  \item 예제 계산을 종이에 직접 다시 수행한다.
  \item 정리의 증명 전략만 보고 증명의 다음 단계를 예상한다.
  \item 완전한 증명을 읽은 뒤 책을 덮고 빈 종이에서 다시 구성한다.
  \item A단계 문제로 정의와 계산을 확인하고, B단계 문제로 핵심 증명을 훈련한다.
  \item C단계 문제는 여러 절을 연결하므로 장 전체를 읽은 뒤 도전한다.
\end{enumerate}

별표가 붙은 문제는 장의 핵심을 대표하거나 뒤의 내용과 직접 연결되는 문제다. 처음 공부할 때에는 별표 문제를 우선 풀고, 나머지 문제는 복습과 심화 학습에 사용할 수 있다.

\section*{원고의 성격과 검수 상태}

이 프로젝트는 공개 저장소에서 집필되는 오픈 교재다. 특정 교재의 번역본이 아니라, 여러 표준 참고문헌을 연구한 뒤 설명·예제·연습문제와 해설을 독자적으로 구성한다. 현재 Volume 1의 군론, 환론과 다항식론은 첫 번째 수학 검수 패스를 마쳤으며, 문장·표기·조판과 외부 독자 검토를 계속 진행하고 있다. 따라서 이 판은 학습에 사용할 수 있는 통합 초안이지만 아직 정식 출판본은 아니다.

원고, 수정 이력과 오류 제보는 다음 공개 저장소에서 관리한다.

\begin{center}
\url{https://github.com/SengangLemon/open-algebra-kr}
\end{center}

AI 도구는 초안 작성, 계산 재검산, 문장 교정과 LaTeX 조판 보조에 사용된다. 최종 구성과 공개 여부를 결정하고, 수학적 내용의 검토 책임을 지는 사람은 프로젝트 유지관리자다. 오류를 발견한 독자의 제보는 이 교재를 실제로 더 나은 공개 학습 자료로 만드는 중요한 기여가 된다.
